\chapter{Introduction}\label{ch:intro}

\section{Why this reproduction}

Large language models handle many table tasks with no task-specific training: writing SQL for a question,
finding a wrong value in a column, matching two schemas. A model fine-tuned for one task usually does better, but
fine-tuning needs labelled examples, and labelling table tasks by hand is slow and expensive.
Table-Specialist~\cite{xing2024tablellmspecialist} proposes to remove the human from that loop. The model generates
its own training data, and a second, \emph{dual} task decides which generated examples are trustworthy enough to train
on.

The idea is attractive because the check is \emph{execution}, not opinion. Take a SQL query and a Spark-Scala program
written for the same question. If both return the same rows on many random samples of the table, they are probably
both right. Two independently written programs rarely make the same mistake on every sample. If the claim holds, it
gives a cheap way to build task specialists from unlabelled tables. It also raises a practical question that the paper
leaves open: the validator throws most generated data away, so what is lost, and can it be recovered?

This project set out to test the paper's central Llama-3.1-8B NL-to-SQL claim on our own hardware, and then to ask
whether the discarded data can be put to use. The claim being tested was fixed at the start, in the project
\texttt{README.md}:

\begin{keypoint}
On benchmarks never used for training, Table-Specialist beats both the vanilla model and fine-tuning without
validation, by more than the variation between seeds. A tie is reported as a valid finding.
\end{keypoint}

\section{Objectives}

The work was organised around two goals, approved as a plan on 2026-11-10 (\fpath{docs/PLAN_NEXT.md}).
\begin{description}[style=nextline,leftmargin=1.2em]
\item[Goal A: match the paper.] Reproduce the paper's Llama-3.1-8B NL-to-SQL numbers (its Table~3), using the paper's
models and datasets as far as they are available. Extend the reproduction to the paper's other tasks where the data
allows: NL-to-Scala, error detection and NL-to-R.
\item[Goal B: improve on it.] Test whether the data the validator rejects can be recovered or exploited, and compare
against established test-time methods.
\end{description}
Two rules governed both goals. First, an unpublished setting is filled with a documented choice and never tuned until
the score lands near the paper's. Second, repo-faithful scores are always reported, with corrected scores beside them
and never in their place.

\section{What this report contains}

\begin{itemize}
\item \textbf{The original method} (\cref{ch:paper}): problem formulation, Algorithms~1--3, data, training and
evaluation, as far as the paper and its code specify them, and where they do not.
\item \textbf{Our methodology} (\cref{ch:method,ch:setup,ch:impl}): every approach we built, with its motivation,
formalisation, settings, seeds and rounds, kept clearly separate from the paper's own method.
\item \textbf{The experimental history} (\cref{ch:history}): a chronological account of how each result led to the
next decision.
\item \textbf{Complete results} (\cref{ch:results-sql,ch:results-ext,ch:results-other}): per seed, per round, per
test set and per prompt, generated directly from the result files.
\item \textbf{Failed approaches} (\cref{ch:failures}): what was tried, what happened, why, and what we learned.
\item \textbf{Ablations, the comparison with the paper, discussion and limitations}
(\cref{ch:ablations,ch:comparison,ch:discussion,ch:limitations}).
\item \textbf{Appendices}: the full configuration, supplementary tables, significance tests, a run registry, the
incident log, a code inventory, a cross-check of headline values, and reproduction instructions.
\end{itemize}

\section{Summary of findings}

\begin{enumerate}
\item \textbf{Approximate reproduction on Instruct.} Llama-3.1-8B-Instruct, queried through the repository's own chat
call, reproduces the paper's NL-to-SQL Table~3 within about 0.03 per test set (Spider excepted;
\cref{sec:sql-instruct}).
\item \textbf{The base checkpoint does not reproduce the paper's vanilla scores.} This is the checkpoint the README
links, and the gap stays unexplained (\cref{sec:sql-vanilla}).
\item \textbf{The paper's qualitative claims hold on base} (\cref{sec:sql-core}): Table-Specialist beats vanilla and
beats no-validation at round~3. The ordering reverses in round~1, which the paper does not report.
\item \textbf{Repair-TS} is our main method contribution (\cref{sec:ext-repair}). It improves base NL-to-SQL
($+0.0421$) and Instruct NL-to-Scala ($+0.0167$), but not Instruct NL-to-SQL. A same-data control separates data
volume from data content.
\item \textbf{Negative results} are reported in full: DPO on validator labels, base self-repair, rounds on Instruct,
and the BIRD scoring artefact (\cref{ch:failures}).
\end{enumerate}

\section{Conventions}\label{sec:conventions}

\begin{description}[style=nextline,leftmargin=1.2em]
\item[Metrics.] NL-to-SQL, Scala and R use \emph{execution accuracy}: the fraction of test cases whose program
returns the gold result. Error detection uses F1 as defined in \cref{sec:paper-metrics}.
  \begin{itemize}
  \item \ZSiv{} is the mean over WikiSQL, WikiTQ, Spider and BIRD with the zero-shot prompt. It is the only aggregate
  comparable with the paper (Text2Analysis is excluded on both sides).
  \item \Mviii{} is the mean over the 8 cells formed by the 4 test sets and the 2 prompts (zero-shot and few-shot).
  \end{itemize}
\item[Scores.] \emph{Repo-faithful} scores use the authors' comparison code, defects included, and are the headline
numbers. \emph{Corrected} scores fix three comparison defects (\cref{sec:scoring}) and are reported alongside.
\item[Seeds and statistics.] Fine-tuned results are seeds 1, 2 and 3 unless stated otherwise. ``Mean $\pm$ SD'' uses
the sample standard deviation ($n-1$). Vanilla models are run once, greedy and deterministic. Significance is a paired
bootstrap over test cases (\cref{sec:stats}).
\item[Dates.] Dates and times are Indian Standard Time (IST, UTC+5:30), as in the project records.
\item[Provenance markers.] Boxes mark where a statement comes from:
  \begin{paperbox}[title={From the original paper},width=\linewidth]
  stated by the paper or its repository;
  \end{paperbox}
  \begin{oursbox}[width=\linewidth]
  our choice or implementation;
  \end{oursbox}
  \begin{observation}[width=\linewidth]a measured fact;\end{observation}
  \begin{interpretation}[width=\linewidth]our reading of facts;\end{interpretation}
  \begin{hypothesis}[width=\linewidth]an explanation not yet tested.\end{hypothesis}
  Paper figures are from the published version's Table~3~\cite{xing2025emnlp}.
\item[Names.] We write TS for Table-Specialist, NV for the no-validation ablation, and repair-TS for our repair
extension. ``Base'' means Llama-3.1-8B and ``Instruct'' means Llama-3.1-8B-Instruct.
\end{description}
